% De Lege Agraria III (de-lege-agraria-3)
% Source: https://raw.githubusercontent.com/PerseusDL/canonical-latinLit/master/data/phi0474/phi011/phi0474.phi011.perseus-lat2.xml
% Fetched by scripts/fetch_latin.py

\ciceroSection{1} Commodius fecissent tribuni plebis, Quirites, si, quae apud vos de me deferunt, ea coram potius me praesente dixissent; nam et aequitatem vestrae disceptationis et consuetudinem superiorum et ius suae potestatis retinuissent. sed quoniam adhuc praesens certamen contentionemque fugerunt, nunc, si videtur eis, in meam contionem prodeant et, quo provocati a me venire noluerunt, revocati saltem revertantur.

\ciceroSection{2} video quosdam, Quirites, strepitu significare nescio quid et non eosdem voltus quos proxima mea contione praebuerunt in hanc contionem mihi rettulisse. qua re a vobis qui nihil de me credidistis ut eam voluntatem quam semper habuistis erga me retineatis peto; a vobis autem quos leviter immutatos esse sentio parvam exigui temporis usuram bonae de me opinionis postulo, ut eam, si quae dixero vobis probabo, perpetuo retineatis; sin aliter, hoc ipso in loco depositam atque abiectam relinquatis.

\ciceroSection{3} completi sunt animi auresque vestrae, Quirites, me gratificantem Septimiis, Turraniis ceterisque Sullanarum adsignationum possessoribus agrariae legi et commodis vestris obsistere. hoc si qui crediderunt, illud prius crediderint necesse est, hac lege agraria quae promulgata est adimi Sullanos agros vobisque dividi, aut denique minui privatorum possessiones ut in eas vos deducamini. si ostendo non modo non adimi cuiquam glebam de Sullanis agris, sed etiam genus id agrorum certo capite legis impudentissime confirmari atque sanciri, si doceo agris eis qui a Sulla sunt dati sic diligenter Rullum sua lege consulere ut facile appareat eam legem non a vestrorum commodorum patrono, sed a Valgi genero esse conscriptam, num quid est causae, Quirites, quin illa criminatione qua in me absentem usus est non solum meam sed etiam vestram diligentiam prudentiamque despexerit?

\ciceroSection{4} caput est legis xl de quo ego consulto, Quirites, neque apud vos ante feci mentionem, ne aut refricare obductam iam rei publicae cicatricem viderer aut aliquid alienissimo tempore novae dissensionis commovere, neque vero nunc ideo disputabo quod hunc statum rei publicae non magno opere defendendum putem, praesertim qui oti et concordiae patronum me in hunc annum populo Romano professus sim, sed ut doceam Rullum posthac in eis saltem tacere rebus in quibus de se et de suis factis taceri velit.

\ciceroSection{5} omnium legum iniquissimam dissimillimamque legis esse arbitror eam quam L. Flaccus interrex de Sulla tulit, ut omnia quaecumque ille fecisset essent rata. nam cum ceteris in civitatibus tyrannis institutis leges omnes exstinguantur atque tollantur, hic rei publicae tyrannum lege constituit. est invidiosa lex, sicuti dixi, verum tamen habet excusationem; non enim videtur hominis lex esse, sed temporis.

\ciceroSection{6} quid si est haec multo impudentior? nam Valeria lege Corneliisque legibus eripitur civi, civi datur, coniungitur impudens gratificatio cum acerba iniuria; sed tamen imbibit illis legibus spem non nullam cui ademptum est, aliquem scrupulum cui datum est. Rulli cautio est haec: ' Qvi post C. Marivm Cn. Papirivm consvles. ' quam procul a suspicione fugit, quod eos consules qui adversarii Sullae maxime fuerunt potissimum nominavit! si enim Sullam dictatorem nominasset, perspicuum fore et invidiosum arbitratus est. sed quem vestrum tam tardo ingenio fore putavit cui post eos consules Sullam dictatorem fuisse in mentem venire non posset?

\ciceroSection{7} quid ergo ait Marianus tribunus plebis, qui nos Sullanos in invidiam rapit? Qvi post Marivm et Carbonem consvles agri, aedificia, lacvs, stagna, loca, possessiones —caelum et mare praetermisit, cetera complexus est—' pvblice data adsignata, vendita, concessa svnt '—a quo, Rulle? post Marium et Carbonem consules quis adsignavit, quis dedit, quis concessit praeter Sullam?—' ea omnia eo ivre sint '—quo iure? labefactat videlicet nescio quid. nimium acer, nimium vehemens tribunus plebis Sullana rescindit—' vt qvae optimo ivre privata svnt. ' etiamne meliore quam paterna et avita?

\ciceroSection{8} meliore . at hoc Valeria lex non dicit, Corneliae leges non sanciunt, Sulla ipse non postulat. si isti agri partem aliquam iuris, aliquam similitudinem propriae possessionis, aliquam spem diuturnitatis attingunt, nemo est tam impudens istorum quin agi secum praeclare arbitretur. tu vero, Rulle, quid quaeris? quod habent ut habeant? quis vetat? Vt privatum sit? ita latum est. Vt meliore iure tui soceri fundus Hirpinus sit sive ager Hirpinus—totum enim possidet—quam meus paternus avitusque fundus Arpinas?

\ciceroSection{9} id enim caves. optimo enim iure ea sunt profecto praedia quae optima condicione sunt. Libera meliore iure sunt quam serva; capite hoc omnia quae serviebant non servient. soluta meliore in causa sunt quam obligata; eodem capite subsignata omnia, si modo Sullana sunt, liberantur. immunia commodiore condicione sunt quam illa quae pensitant; ego Tusculanis pro aqua Crabra vectigal pendam, quia mancipio fundum accepi; si a Sulla mihi datus esset, Rulli lege non penderem.

\ciceroSection{10} video vos, Quirites, sicuti res ipsa cogit, commoveri vel legis vel orationis impudentia, legis quae ius melius Sullanis praediis constituat quam paternis, orationis quae eius modi in causa insimulare quemquam audeat rationes Sullae nimium vehementer defendere. at si illa solum sanciret quae a Sulla essent data, tacerem, modo ipse se Sullanum esse confiteretur. sed non modo illis cavet verum etiam aliud quoddam genus donationis inducit; et is qui a me Sullanas possessiones defendi criminatur non eas solum sancit verum ipse novas adsignationes instituit et repentinus Sulla nobis exoritur.

\ciceroSection{11} nam attendite quantas concessiones agrorum hic noster obiurgator uno verbo facere conetur: ' Qvae data, donata, concessa, vendita. ' patior , audio. quid deinde? ' possessa. ' hoc tribunus plebis promulgare ausus est ut, quod quisque post Marium et Carbonem consules possideret, id eo iure teneret quo quod optimo privatum est ? etiamne si vi deiecit, etiamne si clam, si precario venit in possessionem? ergo hac lege ius civile, causae possessionum, praetorum interdicta tollentur.

\ciceroSection{12} non mediocris res neque parvum sub hoc verbo furtum, Quirites, latet. sunt enim multi agri lege Cornelia publicati nec cuiquam adsignati neque venditi qui a paucis hominibus impudentissime possidentur. his cavet, hos defendit, hos privatos facit; hos, inquam, agros quos Sulla nemini dedit Rullus non vobis adsignare volt, sed eis condonare qui possident. causam quaero cur ea quae maiores vobis in Italia, Sicilia, Africa, duabus Hispaniis, Macedonia, Asia reliquerunt venire patiamini, cum ea quae vestra sunt condonari possessoribus eadem lege videatis.

\ciceroSection{13} iam totam legem intellegetis, cum ad paucorum dominationem scripta sit, tum ad Sullanae adsignationis rationes esse accommodatissimam. nam socer huius vir multum bonus est; neque ego nunc de illius bonitate, sed de generi impudentia disputo. ille enim quod habet retinere volt neque se Sullanum esse dissimulat; hic, ut ipse habeat quod non habet, quae dubia sunt per vos sancire volt et, cum plus appetat quam ipse Sulla, quibus rebus resisto, Sullanas res defendere me criminatur.

\ciceroSection{14} ' habet agros non nullos,' inquit, 'socer meus desertos atque longinquos; vendet eos mea lege quanti volet. habet incertos ac nullo iure possessos; confirmabuntur optimo iure. habet publicos; reddam privatos. denique eos fundos quos in agro Casinati optimos fructuosissimosque continuavit, cum usque eo vicinos proscriberet quoad angulos conformando ex multis praediis unam fundi regionem normamque perfecerit, quos nunc cum aliquo metu tenet, sine ulla cura possidebit.'

\ciceroSection{15} et quoniam qua de causa et quorum causa ille hoc promulgarit ostendi, doceat ipse nunc ego quem possessorem defendam, cum agrariae legi resisto. silvam Scantiam vendis; populus Romanus possidet; defendo. Campanum agrum dividis; vos estis in possessione; non cedo. deinde Italiae, Siciliae ceterarumque provinciarum possessiones venalis ac proscriptas hac lege video; vestra sunt praedia, vestrae possessiones; resistam atque repugnabo neque patiar a quoquam populum Romanum de suis possessionibus me consule demoveri, praesertim, Quirites, cum vobis nihil quaeratur.

\ciceroSection{16} hoc enim vos in errore versari diutius non oportet. num quis vestrum ad vim, ad facinus, ad caedem accommodatus est? nemo . atqui ei generi hominum, mihi credite, Campanus ager et praeclara illa Capua servatur; exercitus contra vos, contra libertatem vestram, contra Cn. Pompeium constituitur; contra hanc urbem Capua, contra vos manus hominum audacissimorum, contra Cn. Pompeium x duces comparantur. veniant et coram, quoniam me in vestram contionem vobis flagitantibus evocaverunt, disserant.
